\chapter{Declarations}
\label{ch:declarations}

\section{Use of Generative AI}
\label{sec:declarations-genai}

I acknowledge the use of Claude Code (Anthropic, \url{https://claude.com/claude-code}) and the underlying Claude Opus / Sonnet models throughout the implementation of the tool described in this thesis. Generative AI was used as a coding assistant in three specific roles, each operating under the developer's direct supervision:

\begin{itemize}
  \item \emph{Repetitive adapter and wrapper code.} The pipeline integrates Eclipse JDT's refactoring catalogue, and each refactoring spec (\texttt{RenameMethod}, \texttt{RenameField}, \texttt{RenameClass}, \texttt{RenameLocalVariable}, \texttt{ExtractMethod}, \texttt{InlineMethod}, \texttt{ChangeMethodSignature}, \texttt{ParameterizeVariable}, and several others) required its own narrowly-distinct wrapper to bridge IntelliJ's recorded refactoring events to JDT's API. The wrappers follow the same structural template, differing only in field-set, precondition, and post-application bookkeeping. AI assistance was used to generate the bulk of each subsequent wrapper from the first hand-written one. The design of the shared interface, the validation logic, and the AST-equivalence check were all done by the developer first.
  \item \emph{Boilerplate around serialisation and test scaffolding.} Kotlin data classes for the analysis-report JSON schema and the fixture builders used throughout the test suite were generated with AI assistance once their shape had been designed.
  \item \emph{Local explanations of unfamiliar APIs.} AI was used to explain corners of the IntelliJ Platform plugin SDK, the Eclipse Equinox bootstrap, and the JCEF dashboard host environment, where official documentation was sparse.
\end{itemize}

Every change was reviewed and integrated by the developer; no AI output was committed without manual review. Behavioural correctness was enforced by a comprehensive test suite across the analysis, refactoring, and dashboard modules, which provides an external check on AI-assisted code independent of inspection.

For the report itself, the content, arguments, and findings are my own. Claude (\url{https://claude.ai}) was used for proofreading, structural feedback on drafts, and assistance with \LaTeX{} formatting of tables and figures. I confirm that no content generated by AI has been presented as my own work.

\section{Ethical Considerations}
\label{sec:declarations-ethics}

The project involves a small user study with two human participants (referred to throughout as P1 and P2) and a contrastive agent comparison with no human subjects. Both participants were senior MEng Computing students at Imperial College London who consented to the study being recorded and analysed, with the understanding that no personally identifying information would appear in the report or the released materials. Participant data (event traces, captured edits) is pseudonymised throughout, and the underlying recordings are kept on the project repository under access controls.

Ethical review was conducted with my project supervisor in line with the Department of Computing's procedure for low-risk human-subjects work; no formal ethics approval was required. No medical, sensitive, or third-party-owned data was used at any stage.

\section{Sustainability}
\label{sec:declarations-sustainability}

The project's compute footprint was minimised by two architectural choices that are described in the methodology chapter and worth highlighting here:

\begin{enumerate}
  \item \emph{Phase A / Phase B pipeline split.} The expensive part of the analysis (shadow-repository reconstruction, RefactoringMiner mining, build/test execution at every checkpoint) is run once per recorded session and cached as a \texttt{phase-a.json} dump. All downstream scoring, sensitivity, ablation, and Monte Carlo experiments re-run only Phase B against the cached dumps, taking seconds rather than the hours a full reproduction would cost. The ablation sweep across 45 sessions and 128 weight variants per session completes in roughly two seconds of wall-clock time on a single machine because of this split.
  \item \emph{Deterministic ranking and per-fixture reuse.} Every experiment driver uses a fixed random seed and produces a stable output, so a rerun for a sanity check costs no more compute than the original. Tie-breaking in the divergence-point ranking is itself deterministic, so the experiments produce bit-identical CSVs across reruns when nothing else changed.
\end{enumerate}

All experiments were run on a single local workstation; no GPU or cloud-compute resources were consumed.

\section{Availability of Data and Materials}
\label{sec:declarations-availability}

The project's materials are accessible as follows:

\begin{enumerate}
  \item \textbf{Source code}: \url{https://github.com/EthanHosier/fyp} -- public repository covering the IntelliJ plugin, the analysis pipeline, the JCEF dashboard, and the experiment drivers.
  \item \textbf{Recorded sessions, fixtures, and supporting material}: every recorded session (45 hand-recorded injection sessions, 30 user-study sessions across 5 participants, 48 agent-comparison sessions across 8 agent configurations), every playbook prompt, the fixtures the recordings were made against, and the session-recording protocol live under \texttt{tool/fixtures/} in the repository. The layout and the meaning of each subdirectory are documented in \texttt{tool/fixtures/README.md}. Per-session JSON artefacts are tracked via Git LFS.
  \item \textbf{Reproducibility of Chapter~\ref{ch:results}}: every numerical claim, table, and plot in Chapter~\ref{ch:results} is reproduced by \texttt{tool/fixtures/notebooks/experiments.ipynb}. The notebook reads the per-session \texttt{phase-a.json} files directly and produces the reported aggregates in-process; setup instructions are in \texttt{tool/fixtures/notebooks/README.md}.
  \item \textbf{Participant data}: pseudonymised; no personally identifying information is included in the repository or in this report.
\end{enumerate}
